\documentclass[a4paper, 12pt]{article}
\usepackage{a4wide}
\usepackage{flexiprogram}
\usepackage{wrapfig}
\usepackage[usenames,dvipsnames]{pstricks}
\usepackage{epsfig}
\usepackage{pst-grad} % For gradients
\usepackage{pst-plot} % For axes
\usepackage{pst-node}
\usepackage{url}
% \usepackage[labelfont={bf}]{caption}

\bibliographystyle{plain}

% \title{Android Based Greenhouse Monitoring And Harvesting Using Accurate And Automated Bot Guidance System}
% \author{Devendra Bhave\\(114050004) \and Wasim\\(114050007) \and Meenakshi Verma
% \\(123050014) \and Mukund Lahoti\\(123050018)}
% \date{\today}

\renewcommand{\baselinestretch}{1.3}
\setlength{\parindent}{0in}

\begin{document}
\begin{titlepage}
 \centering
 \includegraphics[width=0.15\textwidth]{./iitblogo}\\[1.5cm]
 \textsc{\LARGE CS 684 Project Report}\\[1.5cm]
 \textsc{\Large Group \#5} \\ [1.5cm]
 \rule{\linewidth}{0.5mm} \\ [0.5cm]
 {\LARGE \textbf{Greenhouse Monitoring And Harvesting Using Accurate And Automated Bot Guidance System}}
 \rule{\linewidth}{0.5mm}
 \\ [1.5cm]

 \large
 \begin{tabular}{cc}
  Devendra Bhave & Mohd Vasimuddin\\
  (114050004) & (114050007)\\
  \\
  Meenakshi Verma & Mukund Lahoti\\
  (123050014) & (123050018)
 \end{tabular}

 \vspace{2cm}
 \today
\end{titlepage}

\tableofcontents
\thispagestyle{empty}
\cleardoublepage

\newpage
\listoffigures
\listoftables
\thispagestyle{empty}
\cleardoublepage

\setcounter{page}{1}
\section{Introduction}
Greenhouse is the building in which plants are grown. Building is covered with various covering materials like
plastic sheets or glass. This allows solar or artificial light to enter the building, but traps the heat inside
the building. Weather conditions inside a greenhouse can be controlled. Our project aims to
harvest crops in such greenhouses completely autonomously. This report discusses design and implementation aspects of
the project. It talks about design principles applied, engineering choices made and risks mitigated. It also 
enumerates difficulties faced during the project development and further possible enhancements.

\subsection{Definitions, Acronyms and Abbreviations}
\begin{itemize}
 \item FireBird: A robot indigenously designed at ERTS laboratory, IIT Bombay. \cite{eyantra}
 \item AVR-libc: Standard C library implementation by AVR Systems
 \item ABGS: Automatic Bot Guidance System
 \item PWM: Pulse Width Modulation
\end{itemize}

\section{Problem Statement}
Project consists of two key entities -- greenhouse and user. A single greenhouse consists of greenhouse building with
plants arranged in aisles and troughs formation, greenhouse controller and the Bot. There could be multiple such 
greenhouses each one equipped with its own greenhouse controller and the Bot. The entity \emph{user} refers to human farmer.
Greenhouse controller autonomously manages greenhouse functions like sunlight control using sun shades, humidity control using
blowers, etc. It has its own set of sensors, actuators and control logic. User merely specifies its parameters 
like amount of sunlight needed for plants, permissible humidity range, etc. Operation of such greenhouse 
controller is out of scope for this project.

Each greenhouse has one Bot to carry out certain farming tasks. Bot consists of FireBird robot with mounted 
wireless network camera and robotic arm. Figure \ref{fig-bot} shows design of the Bot.
User communicates with Bot via suitable technology. Such communication method is assumed to be 
present. We have used ZigBee point-to-point link for communication between Linux machine and FireBird bot.
User assigns work to the Bot. Bot has capability to move anywhere inside the greenhouse, to take photos and 
videos of plants. It may communicate with greenhouse controller, monitor plant growth, harvest crops and 
alert user when necessary.

\begin{figure}
% Generated with LaTeXDraw 2.0.8
% Fri Nov 16 11:59:57 IST 2012
% \usepackage[usenames,dvipsnames]{pstricks}
% \usepackage{epsfig}
% \usepackage{pst-grad} % For gradients
% \usepackage{pst-plot} % For axes
\scalebox{1} % Change this value to rescale the drawing.
{
\begin{pspicture}(0,-4.92)(15.968125,4.9)
\psframe[linewidth=0.04,dimen=outer](3.6,-0.3)(0.0,-2.5)
\psframe[linewidth=0.04,dimen=outer](5.6,0.7)(2.0,-1.5)
\psarc[linewidth=0.04](3.8,0.7){1.8}{0.0}{180.0}
\psarc[linewidth=0.04](1.8,-0.3){1.8}{0.0}{180.0}
\psline[linewidth=0.04cm](3.6,-2.5)(5.6,-1.5)
\psline[linewidth=0.04cm](3.6,-0.3)(5.6,0.7)
\psline[linewidth=0.04cm](1.0,1.3)(3.0,2.3)
\psline[linewidth=0.04cm](0.0,-2.5)(2.0,-1.5)
\psline[linewidth=0.04cm](0.0,-0.3)(2.0,0.7)
\psline[linewidth=0.04cm](1.8,-1.9)(2.6,-1.9)
\psline[linewidth=0.04cm](1.8,-1.9)(1.8,-2.3)
\psline[linewidth=0.04cm](2.6,-1.9)(2.6,-2.3)
\psline[linewidth=0.04cm](1.8,-2.3)(2.6,-2.3)
\psline[linewidth=0.04cm](1.8,-1.9)(2.2,-1.7)
\psline[linewidth=0.04cm](2.2,-1.7)(3.0,-1.7)
\psline[linewidth=0.04cm](3.0,-1.7)(2.6,-1.9)
\psline[linewidth=0.04cm](3.0,-1.7)(3.0,-2.1)
\psline[linewidth=0.04cm](2.6,-2.3)(3.0,-2.1)
\pscircle[linewidth=0.04,linestyle=dashed,dash=0.16cm 0.16cm,dimen=outer](2.4,-2.1){0.8}
\psline[linewidth=0.04cm](7.0,-0.9)(7.6,-0.3)
\psline[linewidth=0.04cm](11.6,-0.3)(11.0,-0.9)
\psline[linewidth=0.04cm](11.6,-0.3)(11.6,-2.1)
\psline[linewidth=0.04cm](11.6,-2.1)(11.0,-2.7)
\pscircle[linewidth=0.04,dimen=outer](9.1,-2.6){1.1}
\psline[linewidth=0.04cm](10.2,-2.7)(11.0,-2.7)
\psline[linewidth=0.04cm](11.0,-2.7)(11.0,-0.9)
\psline[linewidth=0.04cm](11.0,-0.9)(7.0,-0.9)
\psline[linewidth=0.04cm](7.0,-0.9)(7.0,-2.7)
\psline[linewidth=0.04cm](7.0,-2.7)(8.0,-2.7)
\psarc[linewidth=0.04](9.9,-2.0){1.1}{-80.0}{-40.0}
\psline[linewidth=0.04cm](8.2,2.7)(8.2,0.3)
\psline[linewidth=0.04cm](8.2,2.7)(8.8,3.3)
\psline[linewidth=0.04cm](8.8,3.3)(8.8,0.9)
\psline[linewidth=0.04cm](8.8,0.9)(8.2,0.3)
\psline[linewidth=0.04cm](8.2,2.7)(7.8,2.7)
\psline[linewidth=0.04cm](7.8,2.7)(7.8,0.3)
\psline[linewidth=0.04cm](7.8,2.7)(8.4,3.3)
\psline[linewidth=0.04cm](8.4,3.3)(8.8,3.3)
\psline[linewidth=0.04cm](9.8,0.1)(9.8,-0.7)
\psline[linewidth=0.04cm](10.4,0.1)(10.4,-0.7)
\psline[linewidth=0.04cm](9.8,-0.7)(10.4,-0.7)
\psline[linewidth=0.04cm](10.6,-0.1)(10.6,-0.5)
\psline[linewidth=0.04cm](10.6,-0.5)(10.4,-0.7)
\psline[linewidth=0.04cm](10.6,0.1)(10.6,-0.1)
\psline[linewidth=0.04cm](11.6,-0.3)(10.6,-0.3)
\psline[linewidth=0.04cm](9.8,-0.3)(7.6,-0.3)
\psellipse[linewidth=0.04,dimen=outer](8.5,2.5)(0.1,0.2)
\pscircle[linewidth=0.04,linestyle=dashed,dash=0.16cm 0.16cm,dimen=outer](10.9,0.0){4.9}
\psline[linewidth=0.04cm,linestyle=dashed,dash=0.16cm 0.16cm](2.2,-1.3)(8.2,4.1)
\psline[linewidth=0.04cm,linestyle=dashed,dash=0.16cm 0.16cm](2.2,-2.9)(10.2,-4.9)
\usefont{T1}{ptm}{m}{n}
\rput(9.8385935,-1.19){FireBird Bot}
\usefont{T1}{ptm}{m}{n}
\rput(1.8509375,2.41){Greenhouse}
\usefont{T1}{ptm}{m}{n}
\rput(12.072344,4.01){Robotic Arm}
\psline[linewidth=0.04cm,arrowsize=0.05291667cm 2.0,arrowlength=1.4,arrowinset=0.4]{->}(11.8,3.7)(11.8,1.7)
\psline[linewidth=0.04cm](7.6,0.3)(12.0,0.3)
\psline[linewidth=0.04cm](8.8,0.9)(11.4,0.9)
\psline[linewidth=0.04cm](11.8,0.9)(12.4,0.9)
\psline[linewidth=0.04cm](12.4,0.9)(12.0,0.3)
\psline[linewidth=0.04cm](7.6,0.3)(7.8,0.5)
\psline[linewidth=0.04cm](7.6,0.3)(7.6,0.1)
\psline[linewidth=0.04cm](7.6,0.1)(12.0,0.1)
\psline[linewidth=0.04cm](12.0,0.3)(12.0,0.1)
\psline[linewidth=0.04cm](12.4,0.9)(12.4,0.7)
\psline[linewidth=0.04cm](12.4,0.7)(12.0,0.1)
\pspolygon[linewidth=0.04,fillstyle=solid](11.4,0.7)(12.6,2.9)(12.8,2.9)(11.8,0.7)
\pspolygon[linewidth=0.04,fillstyle=solid](11.2,0.5)(12.4,2.7)(12.6,2.7)(11.6,0.5)
\psframe[linewidth=0.04,dimen=outer,fillstyle=solid](13.0,2.9)(12.2,2.7)
\pspolygon[linewidth=0.04](12.6,2.9)(13.6,3.5)(13.0,2.9)
\pspolygon[linewidth=0.04](13.0,2.9)(14.0,3.3)(13.0,2.7)
\usefont{T1}{ptm}{m}{n}
\rput(14.2448435,1.41){Cutters}
\usefont{T1}{ptm}{m}{n}
\rput(9.69625,4.01){IP Camera}
\psline[linewidth=0.04cm,arrowsize=0.05291667cm 2.0,arrowlength=1.4,arrowinset=0.4]{->}(14.0,1.7)(13.4,2.7)
\psline[linewidth=0.04cm,arrowsize=0.05291667cm 2.0,arrowlength=1.4,arrowinset=0.4]{->}(9.6,3.7)(9.0,3.3)
\psframe[linewidth=0.04,dimen=outer,fillstyle=solid](13.6,-1.3)(11.0,-2.7)
\psline[linewidth=0.04cm](14.2,-0.7)(13.6,-1.3)
\psline[linewidth=0.04cm](14.2,-2.1)(13.6,-2.7)
\psline[linewidth=0.04cm](11.6,-0.7)(14.2,-0.7)
\psline[linewidth=0.04cm](14.2,-0.7)(14.2,-2.1)
\usefont{T1}{ptm}{m}{n}
\rput(14.646406,0.61){Fruit Collector}
\psline[linewidth=0.04cm,arrowsize=0.05291667cm 2.0,arrowlength=1.4,arrowinset=0.4]{->}(14.2,0.3)(13.2,-0.7)
\end{pspicture} 
}

\caption{Design of FireBird Bot}
\label{fig-bot}
\end{figure}

\section{Requirements}
\subsection{Hardware Requirements}
\begin{itemize}
 \item FireBird V robot
 \item Linux machine
 \item Mounted wireless network camera with WiFi
 \item WiFi access point with Internet connectivity
 \item Robotic arm with cutter
 \item Two ZigBee cards
\end{itemize}

\subsection{Functional Requirements}
\begin{enumerate}
 \item User remotely connects to desired greenhouse.
 \item User specifies trough number and automatic Bot guidance system maneuvers the Bot to reach given trough.
 \item User receives live streaming video from the Bot.
 \item Bot uses cutter to cut fruits and vegetables from branches which drop into collector below.
 \item User gets battery level indication and energy estimates for each of the above task.
\end{enumerate}

\subsection{Non Functional Requirements}
\begin{enumerate}
 \item Intuitive GUI design
 \item Good quality video streaming
 \item High speed Internet connection
\end{enumerate}

\subsection{Design Constraints}
User cannot control more than one Bot at a time. There is no support for command broadcasting to multiple Bots.
Area of greenhouse is determined by ZigBee range.

\section{Implementation}
Given the degree of complexity involved, we have divided problem statement into three subtasks:
\begin{enumerate}
 \item[Task \#1]: Move the Bot to desired trough with certain degree of accuracy using automatic bot guidance system 
 \item[Task \#2]: Fetch images from network camera and process it to detect fruits and cutter.
 \item[Task \#3]: Control cutter and robotic arm to precisely cut fruits.
\end{enumerate}

\subsection{FireBird Bot System Architecture}
Figure \ref{fig-arch} shows architecture of system running on the Bot. We have followed modular design methodology.
\begin{figure}
 % Generated with LaTeXDraw 2.0.8
% Fri Nov 16 12:30:13 IST 2012
% \usepackage[usenames,dvipsnames]{pstricks}
% \usepackage{epsfig}
% \usepackage{pst-grad} % For gradients
% \usepackage{pst-plot} % For axes
\scalebox{1} % Change this value to rescale the drawing.
{
\begin{pspicture}(0,-2.5)(16.705,2.5)
\usefont{T1}{ptm}{m}{n}
\rput(5.2203126,-1.93){Firebird}
\usefont{T1}{ptm}{m}{n}
\rput(0.77453125,-0.99){ADC}
\usefont{T1}{ptm}{m}{n}
\rput(2.1764061,-0.99){Buzzer}
\usefont{T1}{ptm}{m}{n}
\rput(3.5403125,-0.99){LCD}
\usefont{T1}{ptm}{m}{n}
\rput(4.9064064,-0.99){Motor}
\usefont{T1}{ptm}{m}{n}
\rput(6.3464065,-0.99){Power}
\usefont{T1}{ptm}{m}{n}
\rput(7.6984377,-0.99){Servo}
\usefont{T1}{ptm}{m}{n}
\rput(9.202969,-0.99){Zigbee}
\psframe[linewidth=0.04,dimen=outer](9.8,-0.5)(0.0,-1.5)
\psline[linewidth=0.04cm](8.4,-0.5)(8.4,-1.5)
\psline[linewidth=0.04cm](7.0,-0.5)(7.0,-1.5)
\psline[linewidth=0.04cm](5.6,-0.5)(5.6,-1.5)
\psline[linewidth=0.04cm](4.2,-0.5)(4.2,-1.5)
\psline[linewidth=0.04cm](2.8,-0.5)(2.8,-1.5)
\psline[linewidth=0.04cm](1.4,-0.5)(1.4,-1.5)
\psframe[linewidth=0.04,dimen=outer](9.8,-1.5)(0.0,-2.5)
\psframe[linewidth=0.04,dimen=outer](9.8,0.5)(0.0,-0.5)
\psline[linewidth=0.04cm](5.0,-0.5)(5.0,0.5)
\usefont{T1}{ptm}{m}{n}
\rput(2.6935937,0.01){Assertions}
\usefont{T1}{ptm}{m}{n}
\rput(7.336406,0.01){AVR-libc}
\psframe[linewidth=0.04,dimen=outer](9.8,1.5)(0.0,0.5)
\psline[linewidth=0.04cm](5.0,1.5)(5.0,0.5)
\usefont{T1}{ptm}{m}{n}
\rput(2.5923438,1.01){ROM Filesystem}
\usefont{T1}{ptm}{m}{n}
\rput(7.868125,1.01){Whiteline Follower}
\psframe[linewidth=0.04,dimen=outer](9.8,2.5)(0.0,1.5)
\usefont{T1}{ptm}{m}{n}
\rput(5.262344,2.01){Bot Guidance System}
\usefont{T1}{ptm}{m}{n}
\rput(11.449062,-1.99){Hardware}
\usefont{T1}{ptm}{m}{n}
\rput(13.612657,-0.99){Hardware Abstraction Layer (HAL)}
\usefont{T1}{ptm}{m}{n}
\rput(12.1925,0.01){Software Libraries}
\usefont{T1}{ptm}{m}{n}
\rput(11.259375,1.01){Utilities}
\usefont{T1}{ptm}{m}{n}
\rput(11.556406,2.01){Application}
\end{pspicture} 
}

 \caption{System Architecture of FireBird Bot}
 \label{fig-arch}
\end{figure}

There are multiple layers of modules, each supporting well defined primitives. Role of each module is 
summarized below:
\begin{itemize}
 \item \textbf{Firebird:} This refers to hardware layer of FireBird Bot. Hardware is controlled by changing
  values in command and status registers. Refer to hardware \cite{fbhwmanual} and software \cite{fbswmanual} 
  manuals of FireBird V for further details.
  
  \item \textbf{HAL Layer:} Hardware abstraction layer (HAL) module hides hardware registers. It offers
  typical driver like primitives. \texttt{init<Device>()} API initializes the hardware of that device. Table \ref{table-hal-primitives} lists all HAL primitives and their use.
  
  \item \textbf{Assertions:} Assertion module mimics C assertions. It supports macro 
  \texttt{ASSERT(condition)}. If assertion fails (\emph{i.e.} condition is false), Bot halts, shows
  line number, source file name and failed statement on LCD screen and starts beeping until reset. Assertion
  is quite useful for debugging. Define macro \texttt{NDEBUG} at the start of the file to turn assertions off. This replaces \texttt{ASSERT()}
  statements by empty statements.
  
  \item \textbf{ROM Filesystem:} This module mimics ROM based file system in limited manner. It redirects standard
  C file streams \texttt{stdin} and \texttt{stdout} to ZigBee. Standard I/O functions  like \texttt{printf()}, \texttt{putchar()}, etc.
  send data over ZigBee to remote console. This simplifies data transmission over ZigBee
  as well as helps in debugging Bot code.  
  \texttt{scanf()} and \texttt{getchar()} functions read from ZigBee. This allows to accept arbitrary data from
  remote machine. Bot can accept arena map files at runtime.
  AVR C library does not support files directly. We have provided support for compiled-in files using predefined
  file handles. File handle named \texttt{MAP\_FILE} can be used to read from compiled-in read-only map file using
  standard buffered I/O functions like \texttt{fscanf(MAP\_FILE, ...)}. More detailed discussion about file system
  and format of map file is included in file \texttt{readme.md} in the project source code.
  
  \item \textbf{White Line Follower:} This module supports white line related operations. It recognizes 
  \emph{checkpoints} (Explained in section \ref{section-bgs}). Table \ref{table-wlf-primitives} explains 
  supported operations. \emph{Orientation} gives direction the Bot is facing towards.
  Bot supports only four possible orientations:
  \begin{itemize}
   \item \texttt{EASTWARD}: Along positive X-axis
   \item \texttt{NORTHWARD}: Along positive Y-axis
   \item \texttt{WESTWARD}: Along negative X-axis
   \item \texttt{SOUTHWARD}: Along negative Y-axis
  \end{itemize}

\end{itemize}


\begin{table}
 \centering
 \begin{tabular}{|l|l|}
 \hline
 \multicolumn{1}{|c|}{\textbf{Primitive}} & \multicolumn{1}{|c|}{\textbf{Use}}\\
 \hline \hline
 \texttt{initAdc()} & Initialize ADC hardware\\
 \texttt{getAdcValue(adc\_channel)} & Read ADC value for specified channel\\
 \hline
 \texttt{initBuzzer()} & Initialize buzzer hardware\\
 \texttt{buzzerOn()} & Turn on buzzer\\
 \texttt{buzzerOff()} & Turn off buzzer\\
 \hline
 \texttt{initLcd()} & Initialize LCD hardware\\
 \texttt{lcdHome()} & Place cursor at first column on LCD display\\
 \texttt{lcdClear()} & Clear LCD screen\\
 \texttt{lcdCursor(row, column)} & Place cursor at given row and column on LCD display\\
 \texttt{lcdString(data)} & Write null terminated data on LCD display\\
 \hline
 \texttt{initMotor()} & Initialize DC motor hardware\\
 \texttt{motorDirectionSet(direction)} & Controls direction of DC motors\\
 \texttt{motorVelocitySet(lvel, rvel)} & Set motor velocity using pulse width modulation\\
 \texttt{motorVelocityGet()} & Read motor velocity settings\\
 \texttt{motorLeftPositionEncoder} & Register left positional encoder callback\\
 \multicolumn{1}{|r|}{\texttt{Init(lCallback)}} & \\
 \texttt{motorRightPositionEncoder} & Register right positional encoder callback\\
 \multicolumn{1}{|r|}{\texttt{Init(rCallback)}} & \\
 \texttt{motorLeftPositionEncoder} & Enable/disable left positional encoder interrupt\\
 \multicolumn{1}{|r|}{\texttt{InterruptConfig(state)}} &  \\
 \texttt{motorRightPositionEncoder} & Enable/disable right positional encoder interrupt\\
 \multicolumn{1}{|r|}{\texttt{InterruptConfig(state)}} & \\
 \hline
 \texttt{initPower()} & Initialize power management hardware\\
 \texttt{powerOn(sensor\_group)} & Turns power on for given group of sensors\\
 \texttt{powerOff(sensor\_group)} & Turns power off for given group of sensors\\
 \hline
 \texttt{initServo()} & Initialize servo motor hardware\\
 \texttt{servoSet(motor, angle)} & Sets given angle for servo motor\\
 \texttt{servoFree(motor)} &  Unlocks servo motors\\
 \hline
 \texttt{initZigbee()} & Initialize ZigBee hardware\\
 \texttt{zigbeeSendByte(u8Data, stream)} & Send one byte data over ZigBee\\
 \texttt{zigbeeReceiveByte(stream)} &  Receive one byte data over ZigBee\\
 \hline
 \end{tabular}
 \caption{FireBird HAL Primitives}
 \label{table-hal-primitives}
\end{table}

\begin{table}
 \centering
 \begin{tabular}{|l|l|}
 \hline
 \multicolumn{1}{|c|}{\textbf{Primitive}} & \multicolumn{1}{|c|}{\textbf{Use}}\\
 \hline \hline
 \texttt{initWhiteLineFollower()} & Initialize white line follower module\\
 \texttt{moveForwardFollwingLine} & Moves along whiteline till specified distance\\
 \multicolumn{1}{|r|}{\texttt{ByDistance(distance)}}& (in millimeter) is covered\\
 \texttt{moveForwardFollwingLine} & Moves along whiteline until checkpoint is hit\\
 \multicolumn{1}{|r|}{\texttt{ByCheckpoint()}} & \\
 \texttt{rotateBot(direction, angle)} & Rotates bot in specified direction by given degrees\\ 
 \hline
 \end{tabular}
 \caption{White Line Follower Primitives}
 \label{table-wlf-primitives}
\end{table}

\subsection{Automatic Bot Guidance System} \label{section-bgs}
Problem of maneuvering the Bot to desired location in the Greenhouse is of fundamental in nature. Every
greenhouse related project needs to solve it. We designed checkpoint based completely automatic bot guidance
system (ABGS) which moves the Bot to desired location in the greenhouse accurately. Maximum error in moving the Bot is
bounded by fixed constant and is independent of the time and distance covered.

ABGS is initialized with the map of the greenhouse. ABGS tracks the Bot using its Cartesian co-ordinates in millimeters.
ABGS moves the Bot along white lines. It uses position encoder to estimate current Bot location. Use of position
encoders introduces error in tracking current bot location due to hardware imprecision. Such error in known as 
\emph{location error}. Location error is bounded by using checkpoints. \emph{Checkpoint} is a co-ordinate on 
the map whose location is \emph{accurately} known. Whenever Bot moves over checkpoint it re-synchronizes its
estimate of current location. Figure \ref{fig-automata} shows automaton used for checkpoint synchronization.
\begin{figure}
 % Generated with LaTeXDraw 2.0.8
% Fri Nov 16 16:31:00 IST 2012
% \usepackage[usenames,dvipsnames]{pstricks}
% \usepackage{epsfig}
% \usepackage{pst-grad} % For gradients
% \usepackage{pst-plot} % For axes
\scalebox{1} % Change this value to rescale the drawing.
{
\begin{pspicture}(0,-4.414219)(15.545313,4.414219)
\usefont{T1}{ptm}{m}{n}
\rput(2.8828125,1.4107813){\psframebox[linewidth=0.04]{Outside Checkpoint Range}}
\usefont{T1}{ptm}{m}{n}
\rput(10.997031,2.6107812){Within Checkpoint Range}
\usefont{T1}{ptm}{m}{n}
\rput(12.579687,-1.3892188){Checkpoint Missed}
\usefont{T1}{ptm}{m}{n}
\rput(3.1803124,-1.3892188){Unexpected Checkpoint Error}
\usefont{T1}{ptm}{m}{n}
\rput(7.996719,-1.3892188){At Checkpoint}
\psarc[linewidth=0.04,arrowsize=0.05291667cm 2.0,arrowlength=1.4,arrowinset=0.4]{->}(3.061875,1.6007812){0.7}{0.0}{180.0}
\psframe[linewidth=0.04,dimen=outer](13.361875,3.1007812)(8.561875,2.3007812)
\rput{-45.0}(0.6594968,8.393728){\psarc[linewidth=0.04,arrowsize=0.05291667cm 2.0,arrowlength=1.4,arrowinset=0.4]{->}(10.461875,3.4007812){0.5}{0.0}{270.0}}
\psframe[linewidth=0.04,dimen=outer](14.361875,-0.89921874)(10.761875,-1.6992188)
\psframe[linewidth=0.04,dimen=outer](9.361875,-0.89921874)(6.561875,-1.6992188)
\psframe[linewidth=0.04,dimen=outer](5.961875,-0.89921874)(0.361875,-1.6992188)
\psline[linewidth=0.04cm,arrowsize=0.05291667cm 2.0,arrowlength=1.4,arrowinset=0.4]{->}(4.761875,1.7007812)(8.561875,2.7007813)
\psline[linewidth=0.04cm,arrowsize=0.05291667cm 2.0,arrowlength=1.4,arrowinset=0.4]{->}(10.961875,2.3007812)(13.561875,-0.89921874)
\psline[linewidth=0.04cm,arrowsize=0.05291667cm 2.0,arrowlength=1.4,arrowinset=0.4]{->}(10.961875,2.3007812)(7.961875,-0.89921874)
\psline[linewidth=0.04cm,arrowsize=0.05291667cm 2.0,arrowlength=1.4,arrowinset=0.4]{->}(2.361875,1.1007812)(1.561875,-0.89921874)
\usefont{T1}{ptm}{m}{n}
\rput(1.0407813,-2.5892189){Legend:}
\usefont{T1}{ptm}{m}{n}
\rput(2.4882812,-2.9892187){D: Distance covered yet}
\usefont{T1}{ptm}{m}{n}
\rput(3.5298438,-3.3892188){$R_L$: Checkpoint range lower bound}
\usefont{T1}{ptm}{m}{n}
\rput(3.5798438,-3.7892187){$R_H$: Checkpoint range upper bound}
\usefont{T1}{ptm}{m}{n}
\rput(5.697344,-4.1892185){C: Checkpoint sensing (1 means checkpoint hit, 0 otherwise)}
\usefont{T1}{ptm}{m}{n}
\rput(2.0545313,2.6107812){<$D < R_L$, $C==0$>}
\usefont{T1}{ptm}{m}{n}
\rput(6.204531,2.2107813){\psframebox*[framesep=0, boxsep=false,fillcolor=white] {<$D == R_L$, $C==0$>}}
\usefont{T1}{ptm}{m}{n}
\rput(11.244532,4.210781){<$D <= R_H$, $C==0$>}
\usefont{T1}{ptm}{m}{n}
\rput(13.314531,-0.18921874){\psframebox*[framesep=0, boxsep=false,fillcolor=white] {<$D > R_H$, $C==X$>}}
\usefont{T1}{ptm}{m}{n}
\rput(2.0545313,0.21078125){\psframebox*[framesep=0, boxsep=false,fillcolor=white] {<$D < R_L$, $C==1$>}}
\usefont{T1}{ptm}{m}{n}
\rput(9.444531,1.0107813){\psframebox*[framesep=0, boxsep=false,fillcolor=white] {<$D <= R_H$, $C==1$>}}
\psline[linewidth=0.04cm,arrowsize=0.05291667cm 2.0,arrowlength=1.4,arrowinset=0.4]{->}(3.561875,1.1007812)(6.961875,-0.89921874)
\usefont{T1}{ptm}{m}{n}
\rput(5.804531,-0.38921875){\psframebox*[framesep=0, boxsep=false,fillcolor=white] {<$D == R_L$, $C==1$>}}
\end{pspicture} 
}

 \caption{Automaton for Checkpoint Synchronization}
 \label{fig-automata}
\end{figure}


ABGS uses Floyd-Warshall all source shortest path algorithm over the greenhouse map. It supports 
\texttt{gotoPosition(x, y)} primitive which moves the Bot from current location to location with co-ordinates
\texttt{(x, y)}. Table \ref{table-abgs-primitives} lists primitives supported by ABGS.

\vspace{\baselineskip}
\textbf{Features of ABGS:}
\begin{itemize}
 \item Gives guarantee that \emph{location error shall never exceed known constant bound}
 \item Precomputes all node shortest paths using Floyd-Warshall algorithm
 \item Uses integer only computations for speed
 \item Uses fast integer square roots
 \item Validated using assertions for manually specified loop invariants
\end{itemize}

\begin{table}
 \centering
 \begin{tabular}{|l|l|}
 \hline
 \multicolumn{1}{|c|}{\textbf{Primitive}} & \multicolumn{1}{|c|}{\textbf{Use}}\\
 \hline \hline
 \texttt{initBotGuidanceSystem(fp, map)} & Initialize ABGS with from file handle \texttt{fp}\\
 \texttt{gotoPosition(x, y)} & \texttt{x} and \texttt{y} are destination co-ordinates in millimeters\\
 \texttt{setBotOrientation(orientation)} & Changes bot orientation\\
 \texttt{gotoForward(distance)} & Moves the Bot forward by given distance in\\
  & millimeters\\
 \hline
 \end{tabular}
 \caption{Automatic Bot Guidance System Primitives}
 \label{table-abgs-primitives}
\end{table}

\subsection{Object Detection Using Computer Vision}
We have used C-URL library to fetch images from wireless network camera on Linux machine. For image
processing and object detection, we used OpenCV 2.4. Our primary interest lies in detecting cutter and multiple 
ripened fruits. We pasted blue sticker on cutter and fruits are assumed to have reddish hue. Fruits should
have certain minimum size. Such requirement serves two purposes. Firstly fruits that are ready for harvesting
will be large enough. Secondly, it filters out spuriously detected objects making detection robust.

\subsection{Fruit Cutting}
Cutter fitted on robotic arm can move along both vertical and horizontal planes. Figure \ref{fig-traj} shows
cutting trajectories followed by robotic arm. Bot starts cutting from one edge and moves forward sweeping
each trajectory. Center of fruit and cutter are aligned by moving robotic arm. When alignment matches,
cutter jaws are closed to cut fruit. Fruit drops in collector bin attached to the Bot. 

\begin{figure}
 % Generated with LaTeXDraw 2.0.8
% Fri Nov 16 17:27:31 IST 2012
% \usepackage[usenames,dvipsnames]{pstricks}
% \usepackage{epsfig}
% \usepackage{pst-grad} % For gradients
% \usepackage{pst-plot} % For axes
\scalebox{1} % Change this value to rescale the drawing.
{
\begin{pspicture}(2,-3.3142188)(17.22,3.3142188)
\psline[linewidth=0.04cm](3.6,-0.19921875)(6.2,1.4007813)
\psline[linewidth=0.04cm](14.6,-0.19921875)(17.2,1.4007813)
\psframe[linewidth=0.04,dimen=outer](14.6,-0.19921875)(3.6,-2.3992188)
\psline[linewidth=0.04cm](17.2,1.4007813)(17.2,-0.7992188)
\psline[linewidth=0.04cm](14.6,-2.3992188)(17.2,-0.7992188)
\psline[linewidth=0.04cm](6.2,1.4007813)(17.2,1.4007813)
\psarc[linewidth=0.04,linestyle=dashed,dash=0.16cm 0.16cm](3.0,-0.19921875){3.0}{0.0}{25.0}
\psarc[linewidth=0.04,linestyle=dashed,dash=0.16cm 0.16cm](5.4,-0.19921875){3.0}{0.0}{32.0}
\psarc[linewidth=0.04,linestyle=dashed,dash=0.16cm 0.16cm](7.6,-0.19921875){3.0}{0.0}{32.0}
\psarc[linewidth=0.04,linestyle=dashed,dash=0.16cm 0.16cm](9.8,-0.19921875){3.0}{0.0}{32.0}
\psarc[linewidth=0.04,linestyle=dashed,dash=0.16cm 0.16cm](12.0,-0.19921875){3.0}{5.0}{32.0}
\psline[linewidth=0.04cm,linestyle=dashed,dash=0.16cm 0.16cm](6.0,-0.19921875)(6.0,-2.3992188)
\psline[linewidth=0.04cm,linestyle=dashed,dash=0.16cm 0.16cm](8.4,-0.19921875)(8.4,-2.3992188)
\psline[linewidth=0.04cm,linestyle=dashed,dash=0.16cm 0.16cm](10.6,-0.19921875)(10.6,-2.3992188)
\psline[linewidth=0.04cm,linestyle=dashed,dash=0.16cm 0.16cm](12.8,-0.19921875)(12.8,-2.3992188)
\psline[linewidth=0.04cm,linestyle=dashed,dash=0.16cm 0.16cm](15.0,0.0)(15.0,-1.9992187)
\usefont{T1}{ptm}{m}{n}
\rput(4.0292187,-3.0892189){Trough}
\psline[linewidth=0.04cm,arrowsize=0.05291667cm 2.0,arrowlength=1.4,arrowinset=0.4]{->}(4.8,-2.9992187)(7.6,-2.3992188)
\usefont{T1}{ptm}{m}{n}
\rput(11.754844,3.1107812){Cutting Trajectories}
\psline[linewidth=0.04cm,arrowsize=0.05291667cm 2.0,arrowlength=1.4,arrowinset=0.4]{->}(10.8,2.8007812)(5.0,0.40078124)
\psline[linewidth=0.04cm,arrowsize=0.05291667cm 2.0,arrowlength=1.4,arrowinset=0.4]{->}(11.0,2.8007812)(7.4,0.60078126)
\psline[linewidth=0.04cm,arrowsize=0.05291667cm 2.0,arrowlength=1.4,arrowinset=0.4]{->}(11.2,2.8007812)(9.4,0.80078125)
\psline[linewidth=0.04cm,arrowsize=0.05291667cm 2.0,arrowlength=1.4,arrowinset=0.4]{->}(11.4,2.8007812)(11.4,1.0007813)
\psline[linewidth=0.04cm,arrowsize=0.05291667cm 2.0,arrowlength=1.4,arrowinset=0.4]{->}(11.6,2.8007812)(13.8,0.80078125)
\psline[linewidth=0.04cm,arrowsize=0.05291667cm 2.0,arrowlength=1.4,arrowinset=0.4]{->}(11.8,2.8007812)(15.8,1.0007813)
\end{pspicture} 
}

 \caption{Cutting Trajectories for Single Trough}
 \label{fig-traj}
\end{figure}


\section{Testing Strategy and Data}
 We have developed testsuits for each of the modules Directory \texttt{FireBird/testsuit} contains all test code.
 \subsection{Testing Automatic Bot Guidance System}
 We have followed method of \emph{assertion based validation}. We have used assertions for sanity checks on all
 parameters and for enforcing loop invariant checks. Any logical error in the module would result in failed assertion. 
 Then, we created test map of size 3000 mm X 3000 mm. We simulated \texttt{gotoPosition()} for all possible pairs
 of co-ordinates. Test procedures \texttt{test\_gotoPosition()} and \texttt{test\_gotoPosition2()} exhaustively 
 test all possible combinations.
 
 \subsection{Testing Fruit Cutting}
 Size of fruit cannot exceed size of cutter jaw. We tested fruit sizes ranging from small beads to table tennis ball.
 We fined tuned colour parameters for small sized objects as small objects are difficult to detect.
 
 \subsection{Energy Consumption Statistics}
 We divided operational battery voltage range into three regions. Table \ref{table-battery} shows all significant
 levels and their associated semantics. Voltage values have been determined by repeated experimentation.
 \begin{table}
  \centering
  \begin{tabular}{|c|c|l|}
  \hline 
  \multicolumn{1}{|c|}{\textbf{Battery Voltage}} & \multicolumn{1}{|c|}{\textbf{Battery API value}} & 
  \multicolumn{1}{|c|}{\textbf{Meaning}}\\
  \hline
  \hline
  9.0 V & $> 900$ & Battery is sufficiently charged.\\
  \hline
  8.0 V & $< 800$ & Battery is low\\
  & & Do not start new task.\\
  & & Finish current task.\\
  \hline
  7.5 V & $< 750$ &  Battery is critically low.\\
  & & Abandon current task.\\
  & & Turn off all servo motors\\
  & & Run towards recharge station.\\
  \hline
  \end{tabular}
  \caption{Battery Voltage Levels}
  \label{table-battery}
 \end{table}
 
 Energy consumption by harvesting task is shown in Table \ref{table-energy}.
 \begin{table}
  \centering
  \begin{tabular}{|l|l|l|}
 \hline
 \multicolumn{1}{|c|}{\textbf{Action}} & \multicolumn{1}{|c|}{\textbf{Energy (Watt-Sec)}} & 
 \multicolumn{1}{|c|}{\textbf{Energy (Watt-Hr)}}\\
 \hline
 \hline
 Move Bot along whiteline & 358 per meter & 0.099 per meter\\
 \hline
 Cutter and arm movement & &\\
 + move forward by 10 cm & 340 & 0.094\\
 \hline
 Scan sideways for fruits & 1326 & 0.368\\
 \hline
 Cutting fruit & 1029 per fruit & 0.286 per fruit\\
 \hline
 One trajectory & 2875 & 0.799\\
 \hline
 One trough & &  \\
 (= 5 trajectories)& 14375 & 3.993\\
 \hline
 \end{tabular}
 \caption{Energy Consumption for Harvesting Task}
 \label{table-energy}
 \end{table}
 
 \section{Design Challenges and Open Issues}
 Ensuring reliability and correctness of the project was the major design concern. We aimed for simple but effective
 Bot design. We chose modular design approach with each module providing well defined services to other modules.
 To ascertain about correctness, we used assertions which enforce invariants at critical points in the code. 
 Any logical error would hopefully violate at least one of these assertions. We used white box testing for each module.
 Testing code is provided in \texttt{FireBird/testsuit} directory. 
 Issues and challenges we faced are listed in following section.
 \paragraph{Achieving Location Accuracy}
 ABGS tracks current location of the Bot using wheel encoders. Hardware inaccuracies introduce small error with
 every movement. Errors accumulate over the time to significant proportion. To highlight specific problems,
 positional encoders do not reliably measure movement less than 4 mm. Both DC motors may not run with same speed
 for same value of PWM. There might be small difference in wheel circumference of each wheel. Response delay of
 each motor may be different. As result of these hardware inaccuracies, achieving same results with repeated trials
 is itself a challenge. We used multiple measures to overcome these issues. First we calibrated hardware and
 incorporated that data into distance calculations. We introduced checkpoints to re-synchronise current location.
 Map of the arena is used to track current location and follow path to destination. This way ABGS is always aware
 of next checkpoint it should expect to hit in its path and how long is it from Bot's current location. 
 Location error did not exceed 7\% (\emph{i.e.} error up to 7 cm in while traveling one meter) in our experiments.
 
 \paragraph{Mitigating Wheel Slipping}
 We observed that robotic arm with the cutter at its front end imparts forward torque on the Bot. This causes center of
 gravity of the Bot to move towards front wheel. As a result, rolling friction between rear wheels and the ground
 reduces and rear wheels slip. Such slipping is undesirable as distance covered by the Bot no longer depends on
 the wheel rotations. Friction between tires and the ground can be increased by either by adding counter weights on rear
 side or by applying better tyre grip. Either approaches increase power consumption of the Bot. We mitigated wheel
 slipping by raising robotic arm to adjust center of gravity before moving forward.
 
 \paragraph{Flow and Error Control for ZigBee}
 ZigBee, being wireless communication medium, is lossy. We designed mechanism to detect errors and recover from them.
 We need flow control mechanism so that Bot receives commands at the rate which it can process. We designed our
 own protocol to meet these requirements. Data is sent over ZigBee using `$!$' terminated fixed sized data packets.
 Receiver replies with either positive or negative acknowledgement. Remember you cannot send raw binary data over
 ZigBee. You must first encode it into printable ASCII characters. Non-printable ASCII characters are reserved
 for serial link control commands. Although this approach serves our purpose, it has lower efficiency. Future
 projects should consider porting TCP/IP protocol stack on FireBird platform.
 
 \paragraph{Map Specification}
 ABGS needs description of the map. Map is supplied to ABGS in predefined format (explained in the code documentation).
 Map file is description of the graph annotated with geographical information for each node. Challenge was how to
 supply map file to ABGS as AVR C library does not suppot files. Thus we developed limited read-only file system
 to meet our purpose. There are two ways in which map file can be read. If map is known beforehand, map file can 
 be compiled into hex code during development and accessed using pre-defined file handle \texttt{MAP\_FILE}.
 Other way to access map file supports runtime loading of the map. We have redirected standard file streams to 
 ZigBee. You can send map file from remote machine over ZigBee. Call map loading API \texttt{loadMap()} with first
 argument as \texttt{stdin}.

 \paragraph{Power Management}
 To minimize power requirement and extend battery time, we added support to turn off unused sensors. 
 \textbf{Instructions from Firebird V software manual \cite{fbswmanual} do not work as intended for sensor power 
 management.} We successfully turned off sharp IR range sensors 1 and 5.
 
 \paragraph{Fruit Detection and Cutting}
 Reliable fruit and cutter detection using OpenCV 2.4 was challenging. We had to pre-process image and fine tune
 colour parameters. Cutting needs correct positioning of the cutter and fruit. Cutting proceeds in two phases. In
 first phase, cutter is positioned approximately near the fruit. In second phase, cutter's position is fine tuned further
 along X and Y axis before cutting. Harvesting task runs as Linux application and sends commands to the Bot using
 ZigBee.
 
\section{Future Work}
We have identified following possible improvements for future projects.
\paragraph{Automatic Recovery of ABGS} Current ABGS can only move Bot along white lines. This is primarily
because ABGS loses Bot's position information if Bot deviates away from white line. Next challenge is drop white line
requirement. Design method which would automatically recover Bot's current location when Bot is not on the white line.

\paragraph{Fruit Depth Information} Single camera cannot give depth information of the fruit to be cut. Thus,
we had to make conservative assumptions about fruit depth. Neither proximity nor range sensors are useful here as 
they fail to distinguish between fruit and leaves while gathering depth information.

\paragraph{Improve Arm Movements} Current cutter design uses single arm. Fruit cutting throughput can be improved 
with complex robotic arm.

\paragraph{Improve Harvesting} Our project is useful for harvesting small fruits like tomatoes, cherries or even
flowers like roses. Future projects should considering harvesting for fruits like grapes, bananas and crops like
corn, sugarcane, etc. Green colour of grapes make their harvesting challenging as detection is difficult.


\section{Conclusion}
This project demonstrates that it is possible to automate harvesting and monitoring in greenhouse. Although real
environment of the greenhouse is rather different from the demonstration platform, this project should still 
function correctly. Choice of right abstractions make this possible. We substantiate this claim for white line. In real
greenhouse, we cannot use white lines because of the soil and the dirt on the floor. However, white line for this project is
mere abstraction of any mechanism that allows the Bot to move in straight line. We can use gyroscope or even
laser guidance system to move along straight line in real greenhouse. This change merely needs modification at HAL
layer code. Rest of the system remains the same.

Embedded systems suffer curse of hardware inaccuracies and physical faults. But, good design can overcome them.
Predictability and accuracy can still be guaranteed within certain known bounds. Layered, module based design and right
abstractions simplify development of complex systems. 


\addcontentsline{toc}{section}{References}
\bibliography{sample}
\end{document}

